\documentclass[11pt,a4paper]{article}
\usepackage[margin=1in]{geometry}
\usepackage{amsmath,amssymb}
\usepackage{booktabs}
\usepackage{array}
\usepackage{hyperref}
\usepackage{xcolor}
\usepackage{listings}
\usepackage{parskip}
\usepackage{titlesec}
\usepackage{enumitem}

\hypersetup{
    colorlinks=true,
    linkcolor=blue!70!black,
    urlcolor=blue!70!black,
}

\lstset{
    basicstyle=\ttfamily\small,
    backgroundcolor=\color{gray!10},
    frame=single,
    breaklines=true,
    columns=fullflexible,
}

\title{\textbf{IBI Graph Pretraining: Design and Analysis}\\
\large Inter-Beat Interval Representation, Self-Supervised Learning,\\and Comparison with wav2vec2/HuBERT}
\author{ECG Foundation Model Project}
\date{\today}

\begin{document}
\maketitle
\tableofcontents
\newpage

% ─────────────────────────────────────────────────────────────────
\section{How Can an IBI Signal Be Represented as a Graph?}
% ─────────────────────────────────────────────────────────────────

An Inter-Beat Interval (IBI) sequence is a time series of durations between consecutive R-peaks in an ECG signal.
Rather than treating it as a flat 1-D sequence, we model it as a graph where each heartbeat is a node and
relationships between beats are encoded as edges.

\subsection{Nodes --- Each Beat Is a Node}

Given a window of $N$ heartbeats, we construct $N$ nodes. Each node $v_i$ carries a feature vector
$\mathbf{x}_i \in \mathbb{R}^{10}$ of hand-crafted, physiologically meaningful attributes.
All features are \emph{MAD-normalized}: $(x - \text{median}) / \text{MAD}$, making them robust to
outliers and invariant across subjects.

\begin{table}[h!]
\centering
\renewcommand{\arraystretch}{1.3}
\begin{tabular}{clll}
\toprule
\textbf{Index} & \textbf{Feature} & \textbf{Description} & \textbf{Clinical Relevance} \\
\midrule
0 & \texttt{ibi\_norm}       & MAD-normalized IBI duration          & Raw HRV magnitude \\
1 & \texttt{dibi\_norm}      & Beat-to-beat change $\Delta$IBI      & Short-term variability (SDNN proxy) \\
2 & \texttt{local\_var\_norm}& Local $\pm$2-beat variance           & Local rhythm stability \\
3 & \texttt{quality}         & Signal quality flag $\in [0,1]$      & Noise-robustness \\
4 & \texttt{hr\_norm}        & Instantaneous heart rate             & Physiological state \\
5 & \texttt{abs\_dibi\_norm} & $|\Delta\text{IBI}|$                 & RMSSD proxy \\
6 & \texttt{pnn50\_flag}     & 1 if $|\Delta\text{IBI}|>50$\,ms    & Classic pNN50 HRV metric \\
7 & \texttt{rel\_ibi}        & IBI / median IBI                     & Beat relative to subject baseline \\
8 & \texttt{trend\_norm}     & Local slope of IBI sequence          & HR acceleration / deceleration \\
9 & \texttt{position}        & Normalized beat index in window      & Temporal context \\
\bottomrule
\end{tabular}
\caption{Node feature vector $\mathbf{x}_i \in \mathbb{R}^{10}$ for beat $i$.}
\end{table}

\subsection{Edges --- Three Types, Built Dynamically}

After the feature MLP projects each node to a 128-dimensional embedding, the adjacency matrix
$A \in \{0,1\}^{B \times N \times N}$ is constructed per batch using three edge types:

\begin{enumerate}[leftmargin=*]

\item \textbf{K-NN Similarity Edges ($K=6$).}
Cosine similarity is computed pairwise over learned embeddings.
Each node is connected to its top-$K$ most similar beats, regardless of temporal distance:
\[
    e_{ij} = 1 \quad \text{if } j \in \mathrm{top\text{-}K}\bigl(\mathrm{cos}(\mathbf{h}_i,\, \mathbf{h}_j)\bigr)
\]
This captures \emph{physiological similarity} --- two beats with similar HRV context are linked
even if they are far apart in time.

\item \textbf{Temporal Sequential Edges.}
Adjacent beats are always connected bidirectionally ($i \leftrightarrow i{+}1$),
preserving the causal structure of the rhythm.

\item \textbf{Self-Loops.}
Every valid node is connected to itself, allowing it to retain its own features through GNN aggregation.

\end{enumerate}

\textbf{Key design property --- dynamic graph construction.}
The graph is \emph{not} precomputed; it is built fresh every forward pass from the current encoder's embeddings.
As training progresses, the encoder improves, and the graph topology improves with it:
beats belonging to similar physiological states cluster together, and those connections reinforce learning.

\subsection{GNN Message Passing}

Each GNN layer performs degree-normalized neighbor aggregation with a residual connection:
\[
    \mathbf{h}_i^{(\ell+1)} = \mathrm{LayerNorm}\!\left(
        \mathrm{GELU}\!\left(
            W_s\,\mathbf{h}_i^{(\ell)}
            + W_n \frac{1}{d_i}\sum_{j} A_{ij}\,\mathbf{h}_j^{(\ell)}
        \right)
        + \mathbf{h}_i^{(\ell)}
    \right)
\]
where $d_i = \sum_j A_{ij}$ is the degree of node $i$.
Three such layers are stacked, followed by attention pooling to obtain a global window-level representation.

% ─────────────────────────────────────────────────────────────────
\section{How to Convert IBI-Graph to Self-Supervised Learning?}
% ─────────────────────────────────────────────────────────────────

The pretraining objective is a \textbf{hybrid} of two complementary SSL signals:
a \emph{local} masked node reconstruction loss and a \emph{global} contrastive (BYOL) loss.

\subsection{Goal 1: Masked Node Reconstruction (Local, Structural)}

\paragraph{Motivation.}
A model that can predict a missing beat's representation from its graph neighbors must learn to understand
rhythmic context --- what surrounding beats imply about the missing one.

\paragraph{Procedure.}
\begin{enumerate}[leftmargin=*]
    \item \textbf{Teacher (no gradient)}: Clean beat features $\rightarrow$ MLP encoder $\rightarrow$ Temporal Transformer $\rightarrow$ normalize $\rightarrow$ target embeddings $\hat{\mathbf{h}}$.
    \item \textbf{Student}: Same clean features $\rightarrow$ MLP encoder $\rightarrow$ replace 30\% of valid nodes with a learnable mask token $\mathbf{m}$ $\rightarrow$ Temporal Transformer $\rightarrow$ build graph $\rightarrow$ GNN $\rightarrow$ linear decoder $\rightarrow$ reconstructed embeddings $\tilde{\mathbf{h}}$.
    \item \textbf{Loss} computed only on masked positions $\mathcal{M}$:
\end{enumerate}
\[
    \mathcal{L}_{\text{mask}} = \frac{1}{|\mathcal{M}|}\sum_{i \in \mathcal{M}}
    \left[
        \|\tilde{\mathbf{h}}_i - \hat{\mathbf{h}}_i\|^2
        + 0.5\,\bigl(1 - \cos(\tilde{\mathbf{h}}_i,\, \hat{\mathbf{h}}_i)\bigr)
    \right]
\]

\paragraph{What the model must learn.}
A masked beat during atrial fibrillation will have very different graph neighbors than one during sinus rhythm.
The graph structure itself becomes the signal: the model must use topology to infer content.

\subsection{Goal 2: BYOL Contrastive Loss (Global, Semantic)}

\paragraph{Motivation.}
Two segments from the same recording should produce the same high-level representation,
invariant to physiological noise and measurement artifacts.

\paragraph{Augmentations applied independently to two views.}
\begin{itemize}[leftmargin=*]
    \item \textbf{Jitter}: Add Gaussian noise $\mathcal{N}(0, (5\,\text{ms})^2)$ to each IBI --- simulates sensor noise.
    \item \textbf{Dropout}: Randomly remove 10\% of beats (keeping $\geq 4$) --- simulates missed R-peak detections.
    \item \textbf{Scale}: Multiply all IBIs by $U(0.92, 1.08)$ --- simulates HR drift or subject movement.
\end{itemize}

\paragraph{Online/Target architecture (BYOL).}
\[
    \mathcal{L}_{\text{BYOL}} = 2 - 2\,\cos\!\bigl(q(\mathbf{z}_1),\,\mathbf{z}_2'\bigr)
    + 2 - 2\,\cos\!\bigl(q(\mathbf{z}_2),\,\mathbf{z}_1'\bigr)
\]
where $\mathbf{z}$, $q(\mathbf{z})$ are the projector and predictor outputs of the \emph{online} network,
and $\mathbf{z}'$ is the projector output of the \emph{target} network (EMA, no gradient).
The EMA momentum is cosine-annealed from $0.996 \to 0.9999$ over training.

\subsection{Combined Objective}

\[
    \boxed{\mathcal{L}_{\text{total}} = \mathcal{L}_{\text{mask}} + 0.5\,\mathcal{L}_{\text{BYOL}}}
\]

The local objective enforces fine-grained beat-level understanding;
the global objective enforces semantic invariance across windows.
Together they prevent representation collapse (BYOL) while preserving structural detail (masked reconstruction).

\subsection{Full Architecture at a Glance}

\begin{lstlisting}
Input: Beat features [B, N, 10]  +  valid_mask [B, N]  +  node_mask [B, N]

  IBI Encoder (MLP 10 -> 64 -> 128)
        |
  Temporal Transformer (2 layers, 4 heads)
        |
  [MASK 30% of valid nodes with learnable token m]
        |
  Build dynamic graph (K-NN + temporal + self-loops)
        |
  GNN (3 layers, normalized aggregation + residual)
        |
  Linear Decoder -> reconstruct masked embeddings
  Attention Pool  -> global representation [B, 128]
        |
  Projector (MLP) -> [B, 128]
  Predictor (MLP) -> BYOL online branch only
\end{lstlisting}

% ─────────────────────────────────────────────────────────────────
\section{How Does IBI-Graph Differ From wav2vec2 / HuBERT?}
% ─────────────────────────────────────────────────────────────────

\subsection{What wav2vec2 and HuBERT Do}

Both methods were designed for \emph{raw audio waveforms} at 16\,kHz sampling rates:
\begin{itemize}[leftmargin=*]
    \item \textbf{wav2vec2}: Convolutional feature extractor $\to$ quantized codebook $\to$ contrastive prediction of masked frames against negative samples.
    \item \textbf{HuBERT}: Offline $k$-means clustering on MFCC / latent representations $\to$ masked prediction of discrete cluster IDs as pseudo-labels.
\end{itemize}
Both treat the signal as a \textbf{dense, fixed-rate 1-D sequence}, and their SSL objectives are
defined in terms of predicting quantized tokens or contrastive frame embeddings.

\subsection{Comparison Table}

\begin{table}[h!]
\centering
\renewcommand{\arraystretch}{1.4}
\begin{tabular}{>{\bfseries}p{3.2cm} p{5.5cm} p{5.5cm}}
\toprule
Dimension & wav2vec2 / HuBERT & IBI-Graph (ours) \\
\midrule
Input modality
    & Dense 1-D waveform ($\sim$16{,}000 samples/sec)
    & Sparse event sequence (0.5--3 beats/sec) \\
Signal structure
    & Short-range temporal correlations
    & Multi-scale: beat-to-beat (ms) + rhythm patterns (minutes) \\
Noise model
    & Additive acoustic noise
    & Ectopic beats, missed detections, variable density \\
Domain knowledge
    & None --- raw signal is target
    & 10 HRV features encode decades of cardiology \\
Cross-subject invariance
    & Waveform morphology varies by lead, body habitus, device
    & MAD-normalized IBI features are subject-invariant \\
Reconstruction target
    & Discrete codebook tokens or quantized frames
    & Continuous teacher embedding (no codebook needed) \\
Relationships captured
    & Strictly sequential (1-D attention)
    & Sequential \textbf{+} similarity-based (graph edges) \\
Variable-rate input
    & Requires fixed-rate resampling
    & Naturally handles irregular beat timing \\
Missing data
    & No explicit mechanism
    & Quality flag node feature + masking \\
\bottomrule
\end{tabular}
\caption{Comparison of IBI-Graph pretraining vs.\ waveform-based SSL methods.}
\end{table}

\subsection{Where IBI-Graph Is a Better Solution}

\paragraph{1. Non-local rhythm patterns via graph topology.}
In wav2vec2 and HuBERT, long-range dependencies must be learned entirely through self-attention from scratch.
In IBI-Graph, K-NN edges \emph{explicitly} connect beats with similar HRV signatures regardless of their
temporal distance. A premature ventricular contraction at beat 3 and another at beat 40 are directly linked
in the graph, making pattern recognition structurally easier.

\paragraph{2. Asynchronous and irregular sequences.}
ECG IBI sequences are fundamentally asynchronous --- beats do not occur at fixed intervals.
wav2vec2 assumes a fixed-rate frame structure. The graph representation handles variable beat timing,
missing beats, and sequence length variation without resampling artifacts.

\paragraph{3. Physiologically grounded features vs.\ learned quantization.}
HuBERT's codebook is derived from $k$-means on raw waveform embeddings --- the resulting clusters
may not align with clinically meaningful categories. The 10-feature node vectors explicitly encode
medically validated HRV metrics (pNN50, RMSSD proxy, local trend), giving the GNN
semantically meaningful structure from the very first training step.

\paragraph{4. Graph topology as an emergent representation.}
As the encoder trains, the K-NN graph that gets built improves.
Beats during sympathetic activation cluster with other sympathetic beats;
parasympathetic recovery with recovery beats.
The \emph{topology itself becomes a learned representation of physiological state} ---
wav2vec2 / HuBERT have no equivalent mechanism.

\paragraph{5. Explicit robustness to signal quality.}
The \texttt{quality} node feature and masked-node mechanism explicitly handle bad beats.
wav2vec2 applied naively to ECG waveforms would be corrupted by motion artifacts and baseline wander
with no selective ignore mechanism.

\subsection{Where IBI-Graph Has Limitations}

\begin{itemize}[leftmargin=*]
    \item \textbf{Lossy abstraction}: IBI discards all waveform morphology (P-wave shape, ST-segment, QRS width) --- useful for rhythm disorders but not for myocardial conditions (e.g., ischemia, hypertrophy).
    \item \textbf{Beat detector dependency}: The entire pipeline fails or degrades if the upstream R-peak detector is inaccurate. wav2vec2 on raw ECG has no such upstream dependency.
    \item \textbf{Shorter effective context}: 30-second windows of $\leq75$ beats vs.\ wav2vec2's ability to process longer audio contexts with hierarchical attention.
\end{itemize}

\subsection{Core Claim}

\begin{center}
\fbox{\parbox{0.85\textwidth}{
\emph{Cardiac rhythm is better modeled as a \textbf{graph of physiologically-featurized events}
than as a waveform sequence, because beats have natural similarity relationships that transcend
temporal ordering, and those relationships carry the information most relevant to autonomic function
and cardiovascular disease state.}
}}
\end{center}

This is a defensible and novel position --- especially for tasks like arrhythmia detection,
stress prediction, sleep staging, and autonomic assessment where HRV is the established clinical gold standard.

\end{document}
